% =============================================================
% CHAPTER 1 — INTRODUCTION
% Purpose: convince the reader the problem matters, state exactly what you did,
% and give them a map of the document. Write a first draft early, then rewrite
% it last so it matches what the thesis actually ends up containing.
% Length: typically 4-8 pages.
% =============================================================
\chapter{Introduction}
\label{chap:introduction}

\section{Context and Motivation}
\guide{Open wide, then narrow. Start from the general phenomenon (deepfakes /
face-swapping, their growth and real-world harms: misinformation, fraud,
non-consensual imagery), supported by a few citations. Explain why detection is
important and why it is hard. Then narrow to your specific angle: using
multimodal LLMs as general-purpose detectors. End the section with the single
sentence that the reader should remember about why this work exists.}

\section{Problem Statement}
\guide{State precisely the problem you tackle, in plain terms. E.g. ``Given a
single face image, decide whether the face has been swapped/manipulated.''
Make the scope explicit: images (not video), specific manipulation families,
zero-shot use of pretrained MLLMs (no fine-tuning), etc. Defining what is OUT of
scope is as valuable as defining what is in.}

\section{Research Questions and Objectives}
\guide{List 2-4 concrete, answerable research questions. These become the
backbone the whole thesis answers. For example:
RQ1: Can off-the-shelf MLLMs detect face swaps above chance in zero-shot mode?
RQ2: How much does prompt design change detection performance?
RQ3: How do results vary across datasets, manipulation generators, and model
families? Then translate each RQ into an objective (what you will build/measure
to answer it).}

\section{Contributions}
\guide{A short bulleted list of what is genuinely new or useful in your work.
Be honest and specific -- a contribution is something a reader could cite. For
this thesis they might be: (i) an extensible, reproducible evaluation pipeline
for MLLM-based manipulation detection; (ii) a systematic comparison of N models
$\times$ M prompts across two datasets; (iii) an analysis of which prompting
strategies and manipulation types help or hurt detection.}
\begin{itemize}
	\item \guide{Contribution 1 -- one sentence.}
	\item \guide{Contribution 2 -- one sentence.}
	\item \guide{Contribution 3 -- one sentence.}
\end{itemize}

\section{Thesis Outline}
\guide{One short paragraph mapping the rest of the document: ``\Cref{chap:background}
introduces the necessary background; \Cref{chap:relatedwork} reviews prior work;
\Cref{chap:methodology} describes the proposed approach; \Cref{chap:setup} details
the experimental setup; \Cref{chap:results} presents and discusses results; and
\Cref{chap:conclusion} concludes and outlines future work.'' Update the
cross-references once your chapters are stable.}
